% ============================================================
%  Проект: Големи јазични модели и формални јазици
%  Предмет: Формални јазици и автомати
%  Студент: Никола Личовски, 196046
%  Наставник: проф. Марија Михова
% ============================================================

\documentclass[12pt,a4paper]{article}

% ── Пакети ──────────────────────────────────────────────────
\usepackage[utf8]{inputenc}
\usepackage[T2A]{fontenc}
\usepackage[macedonian]{babel}
\usepackage{geometry}
\geometry{margin=2.5cm}

\usepackage{amsmath,amssymb,amsthm}
\usepackage{booktabs}
\usepackage{longtable}
\usepackage{array}
\usepackage{multirow}
\usepackage{xcolor}
\usepackage{hyperref}
\usepackage{caption}
\usepackage{subcaption}
\usepackage{float}
\usepackage{enumitem}
\usepackage{listings}
\usepackage{tikz}
\usetikzlibrary{automata, positioning, arrows.meta}
\usepackage{pgfplots}
\pgfplotsset{compat=1.18}
\usepackage{pgfplotstable}
\usepackage{csquotes}
\usepackage{titlesec}
\usepackage{setspace}
% \usepackage[backend=biber, style=numeric, sorting=nyt, language=english]{biblatex}
% \addbibresource{references.bib}
\onehalfspacing

% ── Бои ─────────────────────────────────────────────────────
\definecolor{gptblue}{RGB}{55,138,221}
\definecolor{gemgreen}{RGB}{29,158,117}
\definecolor{codegray}{RGB}{245,245,245}
\definecolor{codecomment}{RGB}{100,100,100}

% ── Стил за код ─────────────────────────────────────────────
\lstset{
  backgroundcolor=\color{codegray},
  basicstyle=\ttfamily\footnotesize,
  commentstyle=\color{codecomment}\itshape,
  keywordstyle=\bfseries,
  breaklines=true,
  frame=single,
  framesep=4pt,
  numbers=left,
  numberstyle=\tiny\color{gray},
  stepnumber=1,
  tabsize=4,
  captionpos=b,
  inputencoding=utf8,
  extendedchars=true,
  literate={а}{{\cyra}}1 {б}{{\cyrb}}1 {в}{{\cyrv}}1
}

% ── Теореми ─────────────────────────────────────────────────
\theoremstyle{definition}
\newtheorem{definition}{Дефиниција}
\newtheorem{theorem}{Теорема}

% ── Хиперврски ──────────────────────────────────────────────
\hypersetup{
  colorlinks=true,
  linkcolor=gptblue,
  citecolor=gptblue,
  urlcolor=gptblue,
  pdftitle={Големи јазични модели и формални јазици},
  pdfauthor={Никола Личовски}
}

% ── Наслов ──────────────────────────────────────────────────
\title{
  \vspace{-1cm}
  {\large Предмет: Формални јазици и автомати}\\[0.8cm]
  \rule{\linewidth}{0.5pt}\\[0.4cm]
  {\LARGE\bfseries Големи јазични модели и формални јазици:\\[0.2cm]
  Препознавање на членство во јазици со модуларни услови}\\[0.4cm]
  \rule{\linewidth}{0.5pt}
}

\author{
  Никола Личовски \quad (196046)\\[0.3cm]
  {\small Наставник: проф. Марија Михова}
}

\date{Јуни 2026}

% ============================================================
\begin{document}
% ============================================================

\maketitle
\thispagestyle{empty}
\newpage

\tableofcontents
\newpage

% ============================================================
\section{Вовед}
% ============================================================

Во последните години, големите јазични модели (Large Language Models --- LLM) покажаа забележителни способности во широк спектар на јазични задачи: генерирање текст, одговарање на прашања, преведување и дури решавање математички проблеми. Меѓутоа, основното прашање --- дали овие модели навистина разбираат и следат строго дефинирани формални структури или само препознаваат статистички шаблони --- останува отворено и е предмет на активно истражување~\cite{deletang2023chomsky,bhattamishra2020ability}.

Формалните јазици претставуваат јазици дефинирани со точни математички правила: граматики, автомати или математички услови. За разлика од природниот јазик, каде контекстот и значењето играат клучна улога, кај формалните јазици припадноста на зборот може да се провери алгоритамски, без двосмисленост. Токму оваа карактеристика ги прави формалните јазици идеален тест за мерење на формалните способности на јазичните модели.

Во овој проект се испитува до кога и под кои услови LLM-моделите успешно препознаваат дали даден збор припаѓа на формално дефиниран јазик. Избраната фамилија јазици се јазиците со \textit{модуларни услови} --- јазици каде бројот на одредени симболи треба да биде делив со даден број. Конкретно, истражувањето анализира три јазика дефинирани врз азбуката $\{a, b, c\}$, каде се брои само симболот $b$, а останатите симболи служат како „шум":

\begin{itemize}
  \item $L_1$: зборови каде бројот на $b$-овци е делив со 3
  \item $L_2$: зборови каде бројот на $b$-овци е делив со 4
  \item $L_3$: зборови каде бројот на $b$-овци е делив со 5
\end{itemize}

Се споредуваат два LLM-модели: ChatGPT (бесплатен интерфејс, претпоставен GPT-4o mini) и Gemini 2.5 Flash (бесплатен интерфејс). Целта е да се измери точноста на секој модел во функција на должина на зборот, апсолутниот број на $b$-овци, распоредот на $b$-овците во зборот, категоријата на зборот (точен, негативен, блиску-точен), класата на остаток и избраниот модул.

Главните истражувачки прашања се:
\begin{enumerate}
  \item До која должина моделите работат сигурно?
  \item При кој апсолутен број на $b$-овци почнуваат систематски да грешат?
  \item Дали распоредот на $b$-овците (групирани наспроти разместени) влијае врз точноста?
  \item Дали постојат систематски разлики меѓу двата модела во насоката на грешките (лажно-позитивни наспроти лажно-негативни)?
  \item Кој модул (3, 4 или 5) е најтежок за препознавање?
  \item Дали моделите користат скратеница базирана на паритет (парно/непарно) наместо вистинско модуларно броење?
\end{enumerate}

\section{Преглед на литература}

Прашањето за формалните способности на трансформерските модели е активно истражувано од неколку аголи.

\vspace{0.8em}

\subsection{Трансформери и регуларни јазици}

\textcite{bhattamishra2020ability} покажаа дека трансформерите теоретски можат да симулираат ограничени автомати, но практичните резултати силно зависат од длабочината и ширината на мрежата, и особено од должината на секвенцата. Нивните наоди покажуваат дека броењето --- операција неопходна за нашиот тип јазици --- е суштински ограничена способност кај моделите без посебни механизми. Подлабока теоретска анализа на тоа кои јазици можат и не можат да ги препознаат трансформерите без надворешна меморија ја дава~\textcite{deletang2023chomsky}, покажувајќи дека строго регуларните јазици теоретски се во дофат на трансформерите, но задачите со прецизно броење ги оптоваруваат во пракса.

\vspace{0.8em}

\subsection{Должинска генерализација}

\textcite{anil2022exploring} го истражуваа проблемот на \textit{length generalization}: дали моделите тренирани на кратки секвенци успешно генерализираат на подолги. Главниот наод е дека LLM-моделите покажуваат значително слабеење на точноста кај должини поголеми од оние видени за време на тренирање, особено кај задачи кои бараат прецизно броење или следење на состојба.

\vspace{0.8em}

\subsection{Модуларна аритметика и LLM}

\textcite{power2022grokking} покажаа дека невралните мрежи можат да научат модуларна аритметика, но само по многу долго тренирање и само кога задачата е формулирана на начин кој им дозволува внатрешно претставување на цикличните структури. Ова сугерира дека LLM-моделите може да имаат внатрешни претставувања за едноставни модуларни операции, но не е јасно дали тоа се пренесува на задача за броење симболи во произволен збор.

\vspace{0.8em}

\subsection{Препознавање на формален јазик преку промпт}

\textcite{hewitt2020rnns} и \textcite{anil2022exploring} истражуваа колку добро ChatGPT и слични модели препознаваат зборови од конкретни формални јазици (на пример $a^nb^n$, правилни загради). Заедничкиот наод е дека моделите покажуваат висока точност за кратки зборови, но брзо деградираат при подолги примери. Исто така, е забележано дека моделите почесто погрешно прифаќаат (лажно-позитивни) отколку погрешно отфрлаат (лажно-негативни) --- наод кој е \textit{во контраст со нашите резултати за ChatGPT}, каде наоѓаме спротивна тенденција.

\vspace{0.8em}

\subsection{Ефект на промпт-инженеринг}

\textcite{wei2022chain} покажуваат дека chain-of-thought промптирање значително ги подобрува резултатите кај аритметички задачи. Во нашето истражување намерно се избегнува chain-of-thought за да се тестира чистата способност на моделот за препознавање на формална структура без водена разработка.

\vspace{0.8em}

\subsection{Резиме}

Постоечката литература конзистентно покажува дека LLM-моделите имаат значителни тешкотии со задачи кои бараат прецизно броење и следење на состојба при долги секвенци. Нашето истражување го проширува ова со систематска анализа на ефектот на апсолутниот број на броени симболи, нивниот просторен распоред, конкретната класа на остаток, и природата на грешките , со споредба меѓу различни модули.

% ============================================================
\section{Формална дефиниција на јазиците}
% ============================================================

\subsection{Азбука и нотација}

Нека $\Sigma = \{a, b, c\}$ е азбуката. За даден збор $w \in \Sigma^*$, со $\#_b(w)$ го означуваме бројот на појавувања на симболот $b$ во $w$.

\subsection{Дефиниција на јазиците}

\begin{definition}[Јазик $L_k$]
За природен број $k \geq 2$, дефинираме:
\[
L_k = \{ w \in \Sigma^* \mid \#_b(w) \equiv 0 \pmod{k} \}
\]
Трите јазика истражувани во овој проект се:
\[
L_3 = \{ w \in \{a,b,c\}^* \mid \#_b(w) \equiv 0 \pmod{3} \}
\]
\[
L_4 = \{ w \in \{a,b,c\}^* \mid \#_b(w) \equiv 0 \pmod{4} \}
\]
\[
L_5 = \{ w \in \{a,b,c\}^* \mid \#_b(w) \equiv 0 \pmod{5} \}
\]
\end{definition}

Симболите $a$ и $c$ не влијаат на членството --- тие служат само како неутрален „шум" кој ја зголемува визуелната комплексност на зборот.

\subsection{Класи на остатоци}

За јазикот $L_k$, секој збор $w \in \Sigma^*$ природно припаѓа на точно една \textit{класа на остатоци} $r \in \{0, 1, \ldots, k-1\}$, каде $r = \#_b(w) \bmod k$. Класата $r=0$ ги опфаќа точно членовите на $L_k$. Оваа структура е клучна за анализата на насоката на грешките:

\begin{itemize}
  \item \textbf{Класа $r=0$ (членови)}: точниот одговор е DA; грешка значи лажно-негативен (ЛН).
  \item \textbf{Класа $r=1$ (блиску-точни)}: еден $b$ далеку од членство; точниот одговор е НЕ; грешка значи лажно-позитивен (ЛП).
  \item \textbf{Класи $r \geq 2$ (негативни)}: подалечни од членство; точниот одговор е НЕ; грешка значи лажно-позитивен (ЛП).
\end{itemize}

\subsection{Примери}

\begin{itemize}
  \item $\varepsilon \in L_3$ бидејќи $\#_b(\varepsilon) = 0 \equiv 0 \pmod{3}$
  \item $\texttt{ababab} \in L_3$ бидејќи $\#_b = 3 \equiv 0 \pmod{3}$
  \item $\texttt{abbca} \notin L_3$ бидејќи $\#_b = 2 \not\equiv 0 \pmod{3}$
  \item $\texttt{bbacbb} \in L_4$ бидејќи $\#_b = 4 \equiv 0 \pmod{4}$
  \item $\texttt{abbbc} \notin L_4$ бидејќи $\#_b = 3 \not\equiv 0 \pmod{4}$
  \item $\texttt{bbcbbab} \in L_5$ бидејќи $\#_b = 5 \equiv 0 \pmod{5}$
  \item $\texttt{abcbc} \notin L_5$ бидејќи $\#_b = 2 \not\equiv 0 \pmod{5}$
\end{itemize}

\subsection{ДКА за секој јазик}

Секој јазик $L_k$ е регуларен и може да се препознае со детерминистички конечен автомат (ДКА) со точно $k$ состојби. Состојбите ја претставуваат тековната вредност $\#_b(w) \bmod k$.

\begin{definition}[ДКА за $L_k$]
Нека $k \geq 2$. Детерминистичкиот конечен автомат кој го препознава
јазикот
\[
L_k=\{\,w\in\{a,b,c\}^* \mid \#_b(w)\equiv 0 \pmod{k}\,\}
\]
е даден со петорката
\[
\mathcal{A}_k=(Q_k,\Sigma,\delta_k,q_0,F_k),
\]
каде:

\begin{itemize}
    \item $Q_k=\{q_0,q_1,\ldots,q_{k-1}\}$ е конечното множество на состојби;
    \item $\Sigma=\{a,b,c\}$ е влезната азбука;
    \item $\delta_k:Q_k\times\Sigma\to Q_k$ е функцијата на премин дефинирана со
    \[
    \delta_k(q_i,b)=q_{(i+1)\bmod k},
    \qquad
    \delta_k(q_i,a)=q_i,
    \qquad
    \delta_k(q_i,c)=q_i;
    \]
    \item $q_0\in Q_k$ е почетната состојба;
    \item $F_k=\{q_0\}$ е множеството на прифатни состојби.
\end{itemize}

Состојбата $q_i$ го претставува остатокот $i=\#_b(w)\bmod k$; по читањето на зборот $w$, автоматот се наоѓа во $q_i$ ако и само ако бројот на досега прочитаните симболи $b$ е конгруентен на $i$ по модул $k$.
\end{definition}


\subsubsection{Дијаграм на состојби за $L_3$}

\begin{figure}[H]
\centering
\begin{tikzpicture}[
  shorten >=2pt,
  node distance=3cm,
  on grid,
  auto,
  >={Stealth[round]}
]
  \node[state, initial, accepting] (q0) {$q_0$};
  \node[state] (q1) [right=of q0] {$q_1$};
  \node[state] (q2) [right=of q1] {$q_2$};

  \path[->]
    (q0) edge [loop above] node {$a,c$} (q0)
    (q0) edge [bend left=20] node {$b$} (q1)
    (q1) edge [loop above] node {$a,c$} (q1)
    (q1) edge [bend left=20] node {$b$} (q2)
    (q2) edge [loop above] node {$a,c$} (q2)
    (q2) edge [bend left=40] node {$b$} (q0);
\end{tikzpicture}
\caption{ДКА $\mathcal{A}_3$ за јазикот $L_3$ (3 состојби)}
\end{figure}

\subsubsection{Дијаграм на состојби за $L_4$}

\begin{figure}[H]
\centering
\begin{tikzpicture}[
  shorten >=2pt,
  node distance=2.8cm,
  on grid,
  auto,
  >={Stealth[round]}
]
  \node[state, initial, accepting] (q0) {$q_0$};
  \node[state] (q1) [right=of q0] {$q_1$};
  \node[state] (q2) [right=of q1] {$q_2$};
  \node[state] (q3) [right=of q2] {$q_3$};

  \path[->]
    (q0) edge [loop above] node {$a,c$} (q0)
    (q0) edge [bend left=20] node {$b$} (q1)
    (q1) edge [loop above] node {$a,c$} (q1)
    (q1) edge [bend left=20] node {$b$} (q2)
    (q2) edge [loop above] node {$a,c$} (q2)
    (q2) edge [bend left=20] node {$b$} (q3)
    (q3) edge [loop above] node {$a,c$} (q3)
    (q3) edge [bend left=40] node[below] {$b$} (q0);
\end{tikzpicture}
\caption{ДКА $\mathcal{A}_4$ за јазикот $L_4$ (4 состојби)}
\end{figure}

\subsubsection{Дијаграм на состојби за $L_5$}

\begin{figure}[H]
\centering
\begin{tikzpicture}[
  shorten >=2pt,
  node distance=2.5cm,
  on grid,
  auto,
  >={Stealth[round]}
]
  \node[state, initial, accepting] (q0) {$q_0$};
  \node[state] (q1) [right=of q0] {$q_1$};
  \node[state] (q2) [right=of q1] {$q_2$};
  \node[state] (q3) [below=of q2] {$q_3$};
  \node[state] (q4) [below=of q0] {$q_4$};

  \path[->]
    (q0) edge [loop above] node {$a,c$} (q0)
    (q0) edge [bend left=20] node {$b$} (q1)
    (q1) edge [loop above] node {$a,c$} (q1)
    (q1) edge [bend left=20] node {$b$} (q2)
    (q2) edge [loop right] node {$a,c$} (q2)
    (q2) edge [bend left=20] node {$b$} (q3)
    (q3) edge [loop below] node {$a,c$} (q3)
    (q3) edge [bend left=20] node {$b$} (q4)
    (q4) edge [loop left] node {$a,c$} (q4)
    (q4) edge [bend left=20] node {$b$} (q0);
\end{tikzpicture}
\caption{ДКА $\mathcal{A}_5$ за јазикот $L_5$ (5 состојби)}
\end{figure}

\subsection{Регуларност на јазиците}

\begin{theorem}
За секое $k \geq 2$, јазикот $L_k$ е регуларен.
\end{theorem}

\begin{proof}
Конструиравме детерминистички конечен автомат
$\mathcal{A}_k=(Q_k,\Sigma,\delta_k,q_0,F_k)$ со $k$ состојби.
По дефиниција, автоматот по читањето на $w$ се наоѓа во $q_i$ ако и само ако $\#_b(w)\equiv i \pmod{k}$.
Оттука, $w$ се прифаќа ако и само ако $\#_b(w)\equiv 0 \pmod{k}$, т.е.\ $w\in L_k$.
Бидејќи секој јазик препознаен со конечен автомат е регуларен, следи дека $L_k$ е регуларен.
\end{proof}

% ============================================================
\section{Алгоритам за точна проверка на членство}
% ============================================================

\subsection{Опис на алгоритмот}

За проверка на членство се користи \textbf{класичен алгоритам за броење симболи} кој го симулира однесувањето на ДКА $\mathcal{A}_k$:

\begin{enumerate}
  \item Го скенира зборот $w$ од лево кон десно, симбол по симбол.
  \item Одржува еден бројач кој ги брои само симболите $b$.
  \item По читање на целиот збор, проверува дали бројачот е делив со $k$.
  \item Враќа \texttt{YES} ако е делив, \texttt{NO} во спротивно.
\end{enumerate}

\subsection{Псевдокод}

\begin{lstlisting}[language=Python, caption={Алгоритам за верификација на членство во $L_k$},
showstringspaces=false,
    showspaces=false]
def verify_membership(word: str, mod: int) -> bool:
    """
    Classical counting algorithm for membership verification.

    Scans the word left to right, counts occurrences of 'b',
    and checks whether that count is divisible by mod.

    Correctness: The language is defined purely by #b(w) mod k.
    This algorithm computes exactly that value with no approximation.

    Complexity: O(n) time, O(1) space

    Args:
        word : the input string over {a, b, c}
        mod  : the modulus k (3, 4, or 5)

    Returns:
        True if #b(word) == 0 (mod k), False otherwise.
    """
    count = sum(1 for ch in word if ch == 'b')
    return count % mod == 0
\end{lstlisting}

\subsection{Коректност и сложеност}

\begin{theorem}
Алгоритмот \texttt{verify\_membership(w, k)} враќа \texttt{True} ако и само ако $w \in L_k$.
\end{theorem}
\begin{proof}
По дефиниција, $w \in L_k \iff \#_b(w) \equiv 0 \pmod{k}$. Алгоритмот ја пресметува вредноста $\#_b(w)$ со еднократно поминување низ $w$ и проверува точно тоа.
\end{proof}

\noindent\textbf{Временска сложеност:} $O(n)$, каде $n = |w|$.
\textbf{Просторна сложеност:} $O(1)$ --- се одржува само еден целоброен бројач.

% ============================================================
\section{Поставување на експериментите}
% ============================================================

\subsection{Генерирање на тест-множеството}

Тест-множеството содржи вкупно \textbf{3.000 зборови} --- по 1.000 за секој јазик ($L_3$, $L_4$, $L_5$). Сите зборови се генерирани програматски со Python скрипта со фиксен случаен seed (\texttt{RANDOM\_SEED = 42}), за целосна репродуцибилност.

\subsubsection{Поделба на категории и класи на остатоци}

Секој јазик содржи три категории зборови:

\begin{table}[H]
\centering
\begin{tabular}{llll}
\toprule
\textbf{Категорија} & \textbf{Број} & \textbf{Класа на остаток} & \textbf{Намена} \\
\midrule
Позитивни   & 400 & $r=0$ & Членови на јазикот \\
Блиску-точни & 200 & $r=1$ & Еден $b$ подалеку од членство \\
Негативни   & 400 & $r \geq 2$ & Не-членови (сите останати остатоци) \\
\bottomrule
\end{tabular}
\caption{Поделба на категории по јазик}
\end{table}

Негативните зборови ги покриваат сите класи $r\in\{2,\ldots,k-1\}$, рамномерно распределени:

\begin{table}[H]
\centering
\begin{tabular}{lll}
\toprule
\textbf{Јазик} & \textbf{Остатоци за негативни} & \textbf{Распределба} \\
\midrule
$L_3$ & $\{2\}$ & 400 зборови \\
$L_4$ & $\{2, 3\}$ & 200 + 200 \\
$L_5$ & $\{2, 3, 4\}$ & 133 + 133 + 134 \\
\bottomrule
\end{tabular}
\caption{Распределба на остатоци за негативните зборови}
\end{table}

\subsubsection{Контролирани варијабли}

\begin{table}[H]
\centering
\begin{tabular}{ll}
\toprule
\textbf{Варијабла} & \textbf{Опис} \\
\midrule
\texttt{language} & Јазикот ($L_3$, $L_4$ или $L_5$) \\
\texttt{length} & Должина на зборот \\
\texttt{residue\_class} & Класа на остаток $r = \#_b(w) \bmod k$ \\
\texttt{b\_count\_tier} & Ниво на апсолутен број $b$-овци \\
\texttt{distribution\_pattern} & Просторен распоред на $b$-овците \\
\bottomrule
\end{tabular}
\end{table}

\subsubsection{Дистрибуција на должини}

\begin{table}[H]
\centering
\begin{tabular}{lrrr}
\toprule
\textbf{Опсег} & \textbf{Дел} & \textbf{Бр.\ зборови (pos/neg)} & \textbf{Бр.\ зборови (near-miss)} \\
\midrule
1 -- 15   & 15\% & 60  & 30  \\
16 -- 30  & 35\% & 140 & 70  \\
31 -- 60  & 35\% & 140 & 70  \\
61 -- 100 & 15\% & 60  & 30  \\
\bottomrule
\end{tabular}
\caption{Дистрибуција на должини на зборови}
\end{table}

\subsubsection{Нивоа на апсолутен број $b$-овци}

\begin{table}[H]
\centering
\begin{tabular}{ll}
\toprule
\textbf{Ниво} & \textbf{Апсолутен број $b$-овци} \\
\midrule
\texttt{very\_low}  & 0 -- 5   \\
\texttt{low}        & 6 -- 15  \\
\texttt{medium}     & 16 -- 30 \\
\texttt{high}       & 31 -- 50 \\
\texttt{very\_high} & 51+      \\
\bottomrule
\end{tabular}
\caption{Нивоа на апсолутен број $b$-овци}
\end{table}

\subsubsection{Просторен распоред на $b$-овците}

Три начини на просторен распоред на $b$-овците се тестирани:

\begin{itemize}
  \item \textbf{Групиран (clustered)}: сите $b$-овци во еден непрекинат блок
  \item \textbf{Наизменичен (interleaved)}: $b$-овци рамномерно распределени
  \item \textbf{Случаен (random)}: позициите избрани случајно (фиксно семе)
\end{itemize}

\subsection{Промпт шаблон}

Ист промпт шаблон е употребен за сите 3.000 тест-примери и за двата модела:

\begin{lstlisting}[caption={Промпт шаблон употребен во сите експерименти}]
You are given a formal language L over the alphabet {a, b, c} defined as:
L = { w | the number of occurrences of 'b' in w is divisible by K }

For each word below, determine whether it belongs to L.

Rules:
- Do not execute any code or use any tools
- Do not explain your reasoning
- Answer with exactly one word per line: YES or NO
- Your response must contain exactly 20 lines in this format:
  N. [word] : YES
  or
  N. [word] : NO

1. [word_1]
2. [word_2]
...
20. [word_20]
\end{lstlisting}

Формалната дефиниција на јазикот е дадена во промптот бидејќи целта е да се тестира способноста за \textit{препознавање на членство}, а не претходното знаење на моделот. Правилото за „без код и алатки" е вклучено бидејќи при почетните тестови беше забележано дека Gemini потенцијално би извршил Python код за броење.

\subsection{Стратегија на испраќање и парсирање}

Зборовите се испраќаат во \textbf{серии од 20 зборови} по промпт (150 серии вкупно). Зборовите се измешани случајно пред батчирањето, за секоја серија да содржи природна мешавина на категории и должини.

При несовпаѓање во одговорот (MISMATCH, MALFORMED, INVALID, MISSING), целата серија се испраќа повторно (максимум 3 обиди). По 3 неуспешни обиди, флагираните зборови се означуваат UNRESOLVABLE и се исклучуваат. Во пракса, нула зборови беа исклучени.

\subsection{Тестирани модели}

\begin{table}[H]
\centering
\begin{tabular}{lll}
\toprule
\textbf{Модел} & \textbf{Интерфејс} & \textbf{Забелешка} \\
\midrule
ChatGPT & Бесплатен, без логирање & Претпоставен GPT-4o mini \\
Gemini  & Бесплатен, без логирање & Gemini 2.5 Flash \\
\bottomrule
\end{tabular}
\caption{Тестирани модели}
\end{table}

Датум на испитување: 15 јуни 2026. Пред секој промпт, страницата беше освежувана за да се избегне влијание на претходните разговори.

\subsection{Ограничувања на методологијата}

\begin{itemize}
  \item \textbf{Неконтролирана температура:} Бесплатните интерфејси не дозволуваат поставување на температура, со што детерминистичноста е ограничена.
  \item \textbf{Неверификувана верзија на моделот:} Точната верзија не може целосно да се потврди.
  \item \textbf{Ограничена репродуцибилност:} Тест-множеството и промптовите се целосно репродуцибилни; одговорите на моделите, меѓутоа, може да варираат при повторно испрашување поради нестохастичност и ажурирања на моделот.
\end{itemize}

% ============================================================
\section{Резултати}
% ============================================================

\subsection{Вкупна точност}

\begin{table}[H]
\centering
\begin{tabular}{lrrrr}
\toprule
\textbf{Модел} & \textbf{Вкупно} & \textbf{Точни} & \textbf{Исклучени} & \textbf{Точност (\%)} \\
\midrule
ChatGPT & 3.000 & 2.434 & 0 & \textbf{81,13} \\
Gemini  & 3.000 & 2.088 & 0 & \textbf{69,60} \\
\bottomrule
\end{tabular}
\caption{Вкупна точност по модел}
\end{table}

ChatGPT го надминува Gemini за \textbf{11,5 процентни поени} на вкупно ниво.

\subsection{Точност по должина на зборот}

\begin{table}[H]
\centering
\begin{tabular}{lrrrrrr}
\toprule
\textbf{Опсег} & \textbf{GPT вкупно} & \textbf{GPT точни} & \textbf{GPT \%} & \textbf{Gem вкупно} & \textbf{Gem точни} & \textbf{Gem \%} \\
\midrule
1 -- 15   & 450  & 444 & 98,67 & 450  & 394 & 87,56 \\
16 -- 30  & 1051 & 898 & 85,44 & 1051 & 769 & 73,17 \\
31 -- 60  & 1051 & 786 & 74,79 & 1051 & 669 & 63,65 \\
61 -- 100 & 448  & 306 & 68,30 & 448  & 256 & 57,14 \\
\bottomrule
\end{tabular}
\caption{Точност по опсег на должина}
\end{table}

\begin{figure}[H]
\centering
\begin{tikzpicture}
\begin{axis}[
  width=12cm, height=6cm,
  xlabel={Опсег на должина},
  ylabel={Точност (\%)},
  symbolic x coords={1--15, 16--30, 31--60, 61--100},
  xtick=data,
  ymin=50, ymax=102,
  legend pos=south west,
  ymajorgrids=true,
  grid style=dashed,
]
\addplot[color=gptblue, mark=square*, thick] coordinates {
  (1--15,98.67) (16--30,85.44) (31--60,74.79) (61--100,68.30)
};
\addplot[color=gemgreen, mark=triangle*, thick, dashed] coordinates {
  (1--15,87.56) (16--30,73.17) (31--60,63.65) (61--100,57.14)
};
\legend{ChatGPT, Gemini}
\end{axis}
\end{tikzpicture}
\caption{Деградација на точноста со зголемување на должината на зборот}
\end{figure}

\subsection{Точност по ниво на апсолутен број $b$-овци}

\begin{table}[H]
\centering
\begin{tabular}{lrrrr}
\toprule
\textbf{Ниво} & \textbf{GPT \%} & \textbf{Gem \%} & \textbf{GPT точни/вк.} & \textbf{Gem точни/вк.} \\
\midrule
very\_low  & 96,48 & 83,79 & 958/993  & 832/993  \\
low        & 84,78 & 66,98 & 724/854  & 572/854  \\
medium     & 70,49 & 64,22 & 461/654  & 420/654  \\
high       & 58,43 & 54,07 & 201/344  & 186/344  \\
very\_high & 58,06 & 50,32 & 90/155   & 78/155   \\
\bottomrule
\end{tabular}
\caption{Точност по ниво на апсолутен број на $b$-овци}
\end{table}

\begin{figure}[H]
\centering
\begin{tikzpicture}
\begin{axis}[
  width=12cm, height=6cm,
  xlabel={Ниво на апсолутен број $b$-овци},
  ylabel={Точност (\%)},
  symbolic x coords={very\_low, low, medium, high, very\_high},
  xtick=data,
  x tick label style={rotate=15, anchor=east},
  ymin=35, ymax=100,
  legend pos=south west,
  ymajorgrids=true,
  grid style=dashed,
]
\addplot[color=gptblue, mark=square*, thick] coordinates {
  (very\_low,96.48) (low,84.78) (medium,70.49) (high,58.43) (very\_high,58.06)
};
\addplot[color=gemgreen, mark=triangle*, thick, dashed] coordinates {
  (very\_low,83.79) (low,66.98) (medium,64.22) (high,54.07) (very\_high,50.32)
};
\legend{ChatGPT, Gemini}
\end{axis}
\end{tikzpicture}
\caption{Точност во функција на апсолутниот број на $b$-овци}
\end{figure}

\subsection{Точност по категорија на зборот}

\begin{table}[H]
\centering
\begin{tabular}{lrrrr}
\toprule
\textbf{Категорија} & \textbf{GPT \%} & \textbf{Gem \%} & \textbf{GPT точни/вк.} & \textbf{Gem точни/вк.} \\
\midrule
Позитивни (positive)      & 72,17 & 69,08 & 866/1200  & 829/1200  \\
Блиску-точни (near\_miss) & 86,83 & 65,00 & 521/600   & 390/600   \\
Негативни (negative)      & 87,25 & 72,42 & 1047/1200 & 869/1200  \\
\bottomrule
\end{tabular}
\caption{Точност по категорија на зборот}
\end{table}

\begin{figure}[H]
\centering
\begin{tikzpicture}
\begin{axis}[
  width=10cm, height=6cm,
  ybar,
  bar width=0.5cm,
  xlabel={Категорија},
  ylabel={Точност (\%)},
  symbolic x coords={positive, near\_miss, negative},
  xtick=data,
  ymin=55, ymax=95,
  legend pos=north west,
  ymajorgrids=true,
  grid style=dashed,
  nodes near coords,
  nodes near coords align={vertical},
  every node near coord/.append style={font=\tiny},
]
\addplot[fill=gptblue, fill opacity=0.85] coordinates {
  (positive,72.17) (near\_miss,86.83) (negative,87.25)
};
\addplot[fill=gemgreen, fill opacity=0.85] coordinates {
  (positive,69.08) (near\_miss,65.00) (negative,72.42)
};
\legend{ChatGPT, Gemini}
\end{axis}
\end{tikzpicture}
\caption{Точност по категорија на зборот}
\end{figure}

\subsection{Точност по јазик (модул)}

\begin{table}[H]
\centering
\begin{tabular}{lrrrr}
\toprule
\textbf{Јазик} & \textbf{GPT \%} & \textbf{Gem \%} & \textbf{GPT точни/вкупно} & \textbf{Gem точни/вкупно} \\
\midrule
$L_3$ (mod 3) & 79,80 & 67,10 & 798/1000 & 671/1000 \\
$L_4$ (mod 4) & 80,60 & 71,00 & 806/1000 & 710/1000 \\
$L_5$ (mod 5) & 83,00 & 70,70 & 830/1000 & 707/1000 \\
\bottomrule
\end{tabular}
\caption{Точност по јазик}
\end{table}

\subsection{Точност по образец на распоред на $b$-овците}

\begin{table}[H]
\centering
\begin{tabular}{lrrrr}
\toprule
\textbf{Образец} & \textbf{GPT \%} & \textbf{Gem \%} & \textbf{GPT точни/вк.} & \textbf{Gem точни/вк.} \\
\midrule
Групиран (clustered)      & 79,50 & 83,58 & 760/956  & 799/956  \\
Наизменичен (interleaved) & 80,06 & 60,55 & 763/953  & 577/953  \\
Случаен (random)          & 83,50 & 65,26 & 911/1091 & 712/1091 \\
\bottomrule
\end{tabular}
\caption{Точност по образец на распоред}
\end{table}

\subsection{Пресечна табела: јазик $\times$ должина}

\begin{table}[H]
\centering
\begin{tabular}{llrrrr}
\toprule
\textbf{Јазик} & \textbf{Опсег} & \textbf{GPT \%} & \textbf{Gem \%} & \textbf{GPT т/в} & \textbf{Gem т/в} \\
\midrule
$L_3$ & 1--15   & 99,33 & 82,67 & 149/150 & 124/150 \\
$L_3$ & 16--30  & 83,14 & 72,29 & 291/350 & 253/350 \\
$L_3$ & 31--60  & 74,86 & 60,57 & 262/350 & 212/350 \\
$L_3$ & 61--100 & 64,00 & 54,67 & 96/150  & 82/150  \\
\midrule
$L_4$ & 1--15   & 97,33 & 89,33 & 146/150 & 134/150 \\
$L_4$ & 16--30  & 84,29 & 72,86 & 295/350 & 255/350 \\
$L_4$ & 31--60  & 73,71 & 65,14 & 258/350 & 228/350 \\
$L_4$ & 61--100 & 71,33 & 62,00 & 107/150 & 93/150  \\
\midrule
$L_5$ & 1--15   & 99,33 & 90,67 & 149/150 & 136/150 \\
$L_5$ & 16--30  & 88,89 & 74,36 & 312/351 & 261/351 \\
$L_5$ & 31--60  & 75,78 & 65,24 & 266/351 & 229/351 \\
$L_5$ & 61--100 & 69,59 & 54,73 & 103/148 & 81/148  \\
\bottomrule
\end{tabular}
\caption{Пресечна табела: јазик $\times$ должина}
\end{table}

\subsection{Пресечна табела: јазик $\times$ категорија}

\begin{table}[H]
\centering
\begin{tabular}{llrrrr}
\toprule
\textbf{Јазик} & \textbf{Категорија} & \textbf{GPT \%} & \textbf{Gem \%} & \textbf{GPT т/в} & \textbf{Gem т/в} \\
\midrule
$L_3$ & positive  & 72,50 & 70,00 & 290/400 & 280/400 \\
$L_3$ & near\_miss & 82,50 & 60,00 & 165/200 & 120/200 \\
$L_3$ & negative  & 85,75 & 67,75 & 343/400 & 271/400 \\
\midrule
$L_4$ & positive  & 69,00 & 66,50 & 276/400 & 266/400 \\
$L_4$ & near\_miss & 90,00 & 68,50 & 180/200 & 137/200 \\
$L_4$ & negative  & 87,50 & 76,75 & 350/400 & 307/400 \\
\midrule
$L_5$ & positive  & 75,00 & 70,75 & 300/400 & 283/400 \\
$L_5$ & near\_miss & 88,00 & 66,50 & 176/200 & 133/200 \\
$L_5$ & negative  & 88,50 & 72,75 & 354/400 & 291/400 \\
\bottomrule
\end{tabular}
\caption{Пресечна табела: јазик $\times$ категорија}
\end{table}

% ── NEW: Residue-class accuracy ────────────────────────────────────────────────
\subsection{Точност по класа на остаток}
\label{subsec:residue_accuracy}

\begin{table}[H]
\centering
\begin{tabular}{lrrrr}
\toprule
\textbf{Класа ($r$)} & \textbf{GPT \%} & \textbf{Gem \%} & \textbf{GPT точни/вк.} & \textbf{Gem точни/вк.} \\
\midrule
$r=0$ (членови)     & 72,17 & 69,08 & 866/1200  & 829/1200  \\
$r=1$ (near-miss)   & 86,83 & 65,00 & 521/600   & 390/600   \\
$r=2$               & 87,60 & 71,39 & 643/734   & 524/734   \\
$r=3$               & 86,49 & 76,58 & 288/333   & 255/333   \\
$r=4$               & 87,22 & 67,67 & 116/133   & 90/133    \\
\bottomrule
\end{tabular}
\caption{Точност по класа на остаток (збирно преку сите јазици)}
\end{table}

Табелата ја открива јасната структурна поделба: класата $r=0$ (вистински членови) е значително потешка за двата модела отколку сите не-членски класи. Кај ChatGPT, точноста по не-членски класи е речиси константна (86--88\%), без оглед на вредноста на остатокот. Кај Gemini, точноста варира значително по не-членски класи ($r=1$: 65,00\% наспроти $r=3$: 76,58\%), со очигледна зависност од „блискоста" до членство.

\begin{table}[H]
\centering
\begin{tabular}{llrrrr}
\toprule
\textbf{Јазик} & \textbf{Класа ($r$)} & \textbf{GPT \%} & \textbf{Gem \%} & \textbf{GPT т/в} & \textbf{Gem т/в} \\
\midrule
$L_3$ & $r=0$ & 72,50 & 70,00 & 290/400 & 280/400 \\
$L_3$ & $r=1$ & 82,50 & 60,00 & 165/200 & 120/200 \\
$L_3$ & $r=2$ & 85,75 & 67,75 & 343/400 & 271/400 \\
\midrule
$L_4$ & $r=0$ & 69,00 & 66,50 & 276/400 & 266/400 \\
$L_4$ & $r=1$ & 90,00 & 68,50 & 180/200 & 137/200 \\
$L_4$ & $r=2$ & 89,50 & 76,50 & 179/200 & 153/200 \\
$L_4$ & $r=3$ & 85,50 & 77,00 & 171/200 & 154/200 \\
\midrule
$L_5$ & $r=0$ & 75,00 & 70,75 & 300/400 & 283/400 \\
$L_5$ & $r=1$ & 88,00 & 66,50 & 176/200 & 133/200 \\
$L_5$ & $r=2$ & 90,30 & 74,63 & 121/134 & 100/134 \\
$L_5$ & $r=3$ & 87,97 & 75,94 & 117/133 & 101/133 \\
$L_5$ & $r=4$ & 87,22 & 67,67 & 116/133 & 90/133  \\
\bottomrule
\end{tabular}
\caption{Точност по јазик $\times$ класа на остаток}
\end{table}

% ── NEW: Error direction tables ────────────────────────────────────────────────
\subsection{Насока на грешките: лажно-негативни и лажно-позитивни}
\label{subsec:error_direction}

Анализата на насоката на грешките разликува два типа:
\begin{itemize}
  \item \textbf{Лажно-негативни (ЛН)}: моделот одговара НЕ за член на јазикот ($r=0$).
  \item \textbf{Лажно-позитивни (ЛП)}: моделот одговара ДА за не-член ($r \geq 1$).
\end{itemize}

\begin{table}[H]
\centering
\begin{tabular}{lrrrr}
\toprule
 & \multicolumn{2}{c}{\textbf{ChatGPT}} & \multicolumn{2}{c}{\textbf{Gemini}} \\
\cmidrule(lr){2-3}\cmidrule(lr){4-5}
\textbf{Јазик} & \textbf{ЛН \%} & \textbf{ЛП \%} & \textbf{ЛН \%} & \textbf{ЛП \%} \\
\midrule
$L_3$ (mod 3) & 27,50 & 15,33 & 30,00 & 34,83 \\
$L_4$ (mod 4) & 31,00 & 11,67 & 33,50 & 26,00 \\
$L_5$ (mod 5) & 25,00 & 11,67 & 29,25 & 29,33 \\
\midrule
\textbf{Вкупно} & \textbf{27,83} & \textbf{12,86} & \textbf{30,92} & \textbf{30,06} \\
\bottomrule
\end{tabular}
\caption{Стапки на лажно-негативни (ЛН) и лажно-позитивни (ЛП) грешки по јазик. ЛН = процент на членови ($r=0$) погрешно отфрлени; ЛП = просечен процент на не-членови погрешно прифатени.}
\label{tab:error_direction}
\end{table}

Табелата го потврдува асиметричниот профил на грешките кај ChatGPT: ЛН стапката (27,83\%) е повеќе од двојно повисока од ЛП стапката (12,86\%), и тоа конзистентно преку сите три јазика. Кај Gemini, ЛН (30,92\%) и ЛП (30,06\%) стапките се речиси еднакви --- симетричен профил на грешки.

Исто така, поделбата по јазик открива важна разлика: ChatGPT има највисока ЛН стапка кај $L_4$ (31,00\%), а најниска кај $L_5$ (25,00\%), со највисок ЛП/ЛН контраст кај $L_4$ (сооднос $\approx$2,7).

\begin{table}[H]
\centering
\begin{tabular}{lrr}
\toprule
\textbf{Јазик (near-miss, $r=1$)} & \textbf{ChatGPT ЛП \%} & \textbf{Gemini ЛП \%} \\
\midrule
$L_3$ ($r=1$) & 17,50 & \textbf{40,00} \\
$L_4$ ($r=1$) & 10,00 & 31,50 \\
$L_5$ ($r=1$) & 12,00 & 33,50 \\
\bottomrule
\end{tabular}
\caption{Лажно-позитивна стапка за блиску-точни зборови ($r=1$) по јазик. Gemini погрешно прифаќа 40\% од near-miss зборовите за $L_3$.}
\label{tab:fp_nearmiss}
\end{table}

\subsection{Консензусни и специфични грешки по модел}

\begin{table}[H]
\centering
\begin{tabular}{lrr}
\toprule
\textbf{Тип на грешка} & \textbf{Број} & \textbf{\% од 3.000} \\
\midrule
И двата модели грешат (консензус)      & 151 & 5,03\% \\
Само ChatGPT греши                      & 415 & 13,83\% \\
Само Gemini греши                       & 761 & 25,37\% \\
И двата модели се точни                 & 1673 & 55,77\% \\
\bottomrule
\end{tabular}
\caption{Распределба на совпаѓања или грешки на моделите}
\end{table}

% ============================================================
\section{Анализа}
% ============================================================

\subsection{Деградација со должина на зборот}

Резултатите јасно потврдуваат дека точноста на двата модела монотоно се намалува со зголемување на должината на зборот. ChatGPT покажува речиси совршена точност за кратки зборови (98,7\% за 1--15), а потоа значително опаѓа при 16--30 (85,4\%), 31--60 (74,8\%) и 61--100 (68,3\%).

\textbf{Критичен праг:} Точката на значителна деградација кај ChatGPT е около должина 16. Кај Gemini, деградацијата е поизразена и при зборови со должина 61--100 достигнува 57,1\% --- граница блиску до случајно погодување (50\% за бинарна задача). Ова е конзистентно со наодите за должинска генерализација~\cite{anil2022exploring}.

\subsection{Апсолутниот број на $b$-овци е посилен предиктор од должината}

Анализата на нивоата на апсолутен број на $b$-овци открива уште поизразена деградација:

\begin{itemize}
  \item ChatGPT: 96,5\% при \texttt{very\_low} → 58,1\% при \texttt{very\_high}
  \item Gemini: 83,8\% при \texttt{very\_low} → 50,3\% при \texttt{very\_high}
\end{itemize}

Gemini при \texttt{very\_high} (51+ $b$-овци) практично броеше случајно. ChatGPT покажа платó --- точноста се стабилизира на околу 58\% веќе од нивото \texttt{high} (31--50 $b$), без значително влошување при \texttt{very\_high}. Ова сугерира дека ChatGPT достигнува горна граница на тешкотија негде помеѓу 30 и 50 $b$-овци.

\textbf{Клучен наод:} Апсолутниот број на симболи кои треба да се бројат е \textit{посилен} предиктор на неуспех отколку вкупната должина на зборот. Долг збор со малку $b$-овци е полесен за моделот отколку краток збор густо исполнет со $b$-овци.

\subsection{Систематска разлика по категорија: предрасуда кон НЕ кај ChatGPT}

ChatGPT е значително полош на позитивни зборови (72,2\%) отколку на негативни (87,3\%) или блиску-точни (86,8\%). Анализата на насоката на грешките (Табела~\ref{tab:error_direction}) покажува дека ЛН стапката (27,83\%) е повеќе од двојно повисока од ЛП стапката (12,86\%). Оваа асиметрија е конзистентна преку сите три јазика, со најизразен контраст кај $L_4$ (ЛН=31,0\%, ЛП=11,67\%, сооднос $\approx$2,7). Се работи за структурна предрасуда, а не случаен шум.

Gemini, за разлика, покажува скоро симетрични стапки (ЛН=30,92\%, ЛП=30,06\%). Ова значи дека Gemini прави речиси ист број грешки кон двете страни --- тоа не е предрасуда кон НЕ туку општа неточност.

\subsection{Структура на грешките по класа на остаток}

Анализата по класа на остаток (Секција~\ref{subsec:residue_accuracy}) открива два различни профила на грешки:

\textbf{ChatGPT}: ЛП стапките по не-членски класи се речиси константни --- 13,17\% ($r=1$), 12,40\% ($r=2$), 13,51\% ($r=3$), 12,78\% ($r=4$), со опсег од само 1,1 процентен поен. Ова значи дека ChatGPT применува скоро иден праг на прифаќање независен од тоа колку е зборот „далеку" од членство.

\textbf{Gemini}: ЛП стапките варираат значително по класа и по јазик. За $L_3$: $r=1$ дава ЛП=40,00\%, $r=2$ дава ЛП=32,25\% --- поблиски зборови до членство се почесто погрешно прифатени. За $L_5$: $r=1$ (ЛП=33,5\%), $r=2$ (ЛП=25,37\%), $r=3$ (ЛП=24,06\%), $r=4$ (ЛП=32,33\%). Интересно е дека $r=4$ за $L_5$ (кое одговара на $r \equiv -1 \pmod 5$, т.е.\ исто така еден $b$ далеку од членство во обратна насока) покажува скок на ЛП слично на $r=1$. Ова сугерира дека Gemini може делумно да ја детектира „блискоста" до членство и во двете насоки по модуларниот циклус.

\subsection{Контраинтуитивниот ефект на модулот}

И двата модела постигуваат \textit{повисока} точност на $L_5$ отколку на $L_3$:

\begin{itemize}
  \item ChatGPT: $L_3$ = 79,8\%, $L_4$ = 80,6\%, $L_5$ = 83,0\%
  \item Gemini: $L_3$ = 67,1\%, $L_4$ = 71,0\%, $L_5$ = 70,7\%
\end{itemize}

Ова е контраинтуитивно бидејќи $L_5$ бара следење на подолг циклус. Исто така, поделбата по ЛН стапка потврдува дека $L_5$ е најлесен и на ниво на грешки ($L_5$: ЛН=25,0\%, наспроти $L_4$: ЛН=31,0\%). Можно објаснување е поголема изложеност на моделите на задачи со делливост со 5 за време на тренирање (делливост со 5 е вообичаена во секојдневната аритметика); сепак, ова е хипотеза и не може да се потврди без увид во тренинг-податоците. Редоследот по тежина е $L_4 > L_3 > L_5$ --- не е монотон во однос на модулот.

\subsection{Одбивање на хипотезата за паритет кај $L_4$}

Можна хипотеза беше дека моделите ја решаваат задачата за $L_4$ преку паритет (парно/непарно) наместо вистинско mod-4 броење: членовите ($r=0$) и $r=2$ се парни, додека $r=1$ и $r=3$ се непарни.

Резултатите од \texttt{mod4\_parity\_hypothesis.csv} ја отфрлаат оваа хипотеза:

\begin{table}[H]
\centering
\begin{tabular}{llllrr}
\toprule
\textbf{Класа} & \textbf{Паритет} & \textbf{Членство} & \textbf{Точен одговор} & \textbf{GPT ДА \%} & \textbf{Gem ДА \%} \\
\midrule
$r=0$ & парно & DA  & ДА & 69,0 & 66,5 \\
$r=1$ & непарно & НЕ & НЕ & 10,0 & 31,5 \\
$r=2$ & парно   & НЕ & НЕ & 10,5 & 23,5 \\
$r=3$ & непарно & НЕ & НЕ & 14,5 & 23,0 \\
\bottomrule
\end{tabular}
\caption{Тест на хипотезата за паритет за $L_4$. Ако моделите користеа паритет, $r=2$ (парно, не-член) би имало значително повисока стапка ДА отколку $r=1$ и $r=3$. Нема такво јаснo разграничување.}
\end{table}

Ако моделите го користеа паритет, непарните не-членови ($r=1$, $r=3$) би добивале НЕ правилно, а парниот не-член ($r=2$) би бил систематски погрешно прифатен. Тоа не се набљудува: кај ChatGPT, $r=2$ (10,5\% ДА) е речиси идентичен со $r=1$ (10,0\% ДА). Кај Gemini, $r=1$ (31,5\%) е дури \textit{повисок} од $r=2$ (23,5\%). Хипотезата за паритет е јасно отфрлена.

\subsection{Асиметрија по начин на распоред --- најизненадувачкиот наод}

Резултатите по образец на распоред откривааат драматична асиметрија:

\begin{itemize}
  \item \textbf{Gemini}: 83,6\% на групирани зборови, само 60,6\% на наизменични --- разлика од \textbf{23 процентни поени}.
  \item \textbf{ChatGPT}: речиси никаков ефект: 79,5\% (групирани), 80,1\% (наизменични), 83,5\% (случајни).
\end{itemize}

Ова е еден од најзначајните наоди. Блокот на $b$-овци (на пример \texttt{aaabbbbbbaaa}) изгледа дека Gemini го обработува полесно отколку кога $b$-овците се расфрлани низ зборот, сугерирајќи стратегија поблиска до \textbf{визуелно препознавање на блок} отколку секвенцијално броење. ChatGPT, пак, покажува конзистентен процес независно од распоредот --- ова е конзистентно со примена на нешто поблиску до секвенцијална обработка на целата влезна секвенца.

\subsection{Споредба на моделите}

ChatGPT конзистентно го надминува Gemini кај сите варијабли освен групираниот образец (каде Gemini е незначително подобар: 83,6\% vs 79,5\%). Разликата е особено изразена кај:

\begin{itemize}
  \item Наизменичен распоред: 80,1\% (GPT) vs 60,6\% (Gem) --- \textbf{19,5 пп}
  \item Блиску-точни зборови ($r=1$): 86,8\% (GPT) vs 65,0\% (Gem) --- \textbf{21,8 пп}
  \item $L_3$ near-miss специфично: 82,5\% (GPT) vs 60,0\% (Gem) --- \textbf{22,5 пп}
\end{itemize}

Квалитативно, двата модела имаат различни профили на неуспех: ChatGPT е систематски скептичен (предрасуда кон НЕ со инваријантна ЛП стапка), додека Gemini е систематски несигурен (симетрични грешки со чувствителност кон блискоста на остатокот до нула).

\subsection{Систематски vs случајни грешки}

Консензус-грешките (и двата модела грешат) сочинуваат само 5,03\% (151 зборови), додека специфичните грешки на Gemini (25,37\%) далеку ги надминуваат оние на ChatGPT (13,83\%). Ова потврдува дека разликата не е случајна --- Gemini систематски греши на зборови кај кои ChatGPT е точен. Консензус-грешките се концентрирани кај \texttt{high}/\texttt{very\_high} ниво на $b$-овци и долги зборови, каде и двата модела достигнуваат граница на способноста за броење.

% ============================================================
\section{Заклучок}
% ============================================================

\subsection{Одговори на истражувачките прашања}

\textbf{1. До која должина моделите работат сигурно?}

ChatGPT работи сигурно за зборови со должина до 15 (98,7\%). Значителна деградација почнува при должина 16. При должини над 60, точноста опаѓа на 68,3\%. Gemini деградира порано и при должини над 60 достигнува 57,1\% --- блиску до случајно погодување.

\textbf{2. При кој апсолутен број на $b$-овци почнуваат да грешат?}

Деградацијата почнува при 6--15 $b$-овци и е најизразена меѓу \texttt{low} и \texttt{medium} (16--30). При повеќе од 30 $b$-овци, двата модела достигнуваат значителна несигурност. Gemini при 51+ $b$-овци практично брои случајно (50,3\%).

\textbf{3. Дали распоредот на $b$-овците влијае?}

Кај ChatGPT, речиси не. Кај Gemini, влијанието е драматично (23 пп разлика), сугерирајќи различни внатрешни стратегии кај двата модела.

\textbf{4. Дали постојат систематски разлики во насоката на грешките?}

Да. ChatGPT покажува силна лажно-негативна предрасуда: ЛН стапката (27,83\%) е повеќе од двојно повисока од ЛП стапката (12,86\%), конзистентно преку сите три јазика. Оваа предрасуда е најизразена кај $L_4$ (ЛН=31,0\%, ЛП=11,67\%, сооднос $\approx$2,7). Gemini покажува речиси симетрични стапки (ЛН$\approx$ЛП$\approx$30\%), но со висока ЛП стапка кај near-miss зборови --- особено 40,0\% за $L_3$.

\textbf{5. Кој модул е најтежок?}

Контраинтуитивно, $L_3$ е најтежок и $L_5$ е најлесен за двата модела. Редоследот по тежина е $L_4 > L_3 > L_5$. Ова не е монотон во однос на модулот и можно објаснување е поголема изложеност на моделите на аритметика со 5 за време на тренирање, иако тоа останува хипотеза.

\textbf{6. Дали моделите користат скратеница со паритет?}

Не. Хипотезата за паритет кај $L_4$ е јасно отфрлена: $r=2$ (парен не-член) не е систематски погрешно прифатен повеќе отколку непарните не-членови.

\subsection{Главни наоди}

\begin{enumerate}
  \item Двата LLM-модели покажуваат јасна монотона деградација на точноста со зголемување на должината, со критичен праг околу должина 15--20.
  \item Апсолутниот број на симболи кои треба да се бројат е \textit{посилен} предиктор на неуспех отколку вкупната должина на зборот.
  \item ChatGPT покажува силна, структурна лажно-негативна предрасуда (ЛН$\approx$2,2$\times$ЛП), со речиси инваријантна ЛП стапка (12--14\%) преку сите не-членски класи на остатоци.
  \item Gemini покажува симетрични стапки на грешки (ЛН$\approx$ЛП$\approx$30\%), но со чувствителност на ЛП стапките кон блискоста на остатокот до нула --- сугерирајќи различен механизам на препознавање.
  \item Gemini е драматично чувствителен на просторниот распоред на $b$-овците (разлика од 23 пп помеѓу групиран и наизменичен), додека ChatGPT не е.
  \item ChatGPT (81,1\%) значително го надминува Gemini (69,6\%), со највисока разлика кај наизменичен распоред (19,5 пп) и блиску-точни зборови (21,8 пп).
  \item Хипотезата за паритет кај $L_4$ е отфрлена: ниту еден модел не прифаќа парните не-членови ($r=2$) значително почесто отколку непарните не-членови.
  \item Ниту еден модел не успева да одржи висока точност при повеќе од 30 $b$-овци, независно од должината.
\end{enumerate}

\subsection{Идни правци}

\begin{itemize}
  \item Тестирање со API-пристап за контролирана температура и потврдени верзии на моделите.
  \item Анализа на chain-of-thought промптирање~\cite{wei2022chain} --- дали образложувањето ја подобрува точноста.
  \item Проширување на јазичната фамилија (на пример, $a^nb^n$, правилно заградени зборови~\cite{deletang2023chomsky}).
  \item Тестирање на поголеми модели (GPT-4o, Gemini Ultra) за споредба.
  \item Детална анализа на конкретните грешки за идентификација на внатрешните стратегии, со посебен акцент на тоа зошто $r=4 \equiv -1 \pmod 5$ предизвикува зголемена ЛП кај Gemini.
\end{itemize}

% ============================================================
\printbibliography[heading=bibintoc, title={Литература}]
% ============================================================

% ============================================================
\appendix
\section{Код}
% ============================================================

Целиот код е достапен во следните Python датотеки:

\begin{description}
  \item[\texttt{1\_generator.py}] Генерирање на тест-множеството (3.000 зборови, семе 42)
  \item[\texttt{2\_prompt\_builder.py}] Конструкција на батч-промптови (150 датотеки)
  \item[\texttt{3\_parser.py}] Парсирање и евидентирање на одговорите
  \item[\texttt{4\_analysis.py}] Пресметување на сите табели и анализи
\end{description}


% ============================================================
\end{document}
